For every order \(\nu\), write \(\lambda_{j,\nu}=\operatorname{cum}_\nu(\varepsilon_j)\). Moment determinacy presupposes finite moments of every order. Mutual independence makes mixed source cumulants vanish, while multilinearity of cumulants gives the observable tensor
\[
 \mathcal T_\nu=\operatorname{cum}_\nu(X)
 =\sum_{j=1}^n\lambda_{j,\nu}c_j^{\otimes\nu}. \tag{50}
\]
Let \(A_\nu=\{j:\lambda_{j,\nu}\ne0\}\) and \(R_\nu=|A_\nu|\). Suppose first that \(R_\nu\ge2\) and \(\nu\ge2R_\nu-1\). Put
\[
 \ell_1=\ell_2=R_\nu-1,
 \qquad
 \ell_3=\nu-2R_\nu+2\ge1. \tag{51}
\]
Reshaping (50) into these three blocks yields
\[
 \widetilde{\mathcal T}_\nu
 =\sum_{j\in A_\nu}\lambda_{j,\nu}
 \operatorname{vec}(c_j^{\otimes\ell_1})\otimes
 \operatorname{vec}(c_j^{\otimes\ell_2})\otimes
 \operatorname{vec}(c_j^{\otimes\ell_3}). \tag{52}
\]
The Veronese--Lagrange rank lemma applies to each of the first two factors because \(\ell_1=\ell_2=R_\nu-1\); their Kruskal ranks equal \(R_\nu\). Any two columns of the third factor are linearly independent: applying the same lemma to that pair uses only \(\ell_3\ge1\). Hence its Kruskal rank is at least two. Therefore
\[
 k_1+k_2+k_3\ge R_\nu+R_\nu+2=2R_\nu+2. \tag{53}
\]
The established Kruskal three-factor uniqueness lemma applied to (52) proves essential uniqueness up to a common permutation and reciprocal nonzero scalings. The tensor rank is exactly \(R_\nu\): the first factor matrix has full column rank, and the Khatri--Rao product of the second and third factors has full column rank. For the latter claim, apply a left-inverse row selecting one column of the second factor to a linear dependence among the Khatri--Rao columns; the selected nonzero third-factor column forces that coefficient to vanish, and repeat for every column. Thus the mode-one matricization has rank \(R_\nu\), furnishing the matching lower bound on tensor rank. If \(R_\nu=1\), a nonzero tensor in (50) is rank one and its common projective preimage is immediate; if \(R_\nu=0\), the tensor is zero and contributes no direction. This proves the adaptive elbow \(\nu\ge2R_\nu-1\).

For a recovery rule that knows \(n\) but not \(R_\nu\), take any \(\nu\ge2n-1\) and use the fixed degrees
\[
 (\ell_1,\ell_2,\ell_3)=(n-1,n-1,\nu-2n+2). \tag{54}
\]
For every active rank \(R_\nu\le n\), the first two degrees are at least \(R_\nu-1\), while the third is at least one. The argument (52)--(53) therefore applies unchanged. A computable exact population operation is the following constrained decomposition: minimize \(R\in\{0,\ldots,n\}\) over nonzero coefficients \(\gamma_1,\ldots,\gamma_R\), nonzero vectors \(z_1,\ldots,z_R\in\mathbb R^p\), and block factors \(b_{s\ell}\), subject to
\[
 \widetilde{\mathcal T}_\nu
 =\sum_{\ell=1}^R\gamma_\ell b_{1\ell}\otimes b_{2\ell}\otimes b_{3\ell},
 \qquad
 b_{s\ell}=\operatorname{vec}(z_\ell^{\otimes\ell_s})
 \quad(s=1,2,3). \tag{55}
\]
Equations (52)--(54), minimality, and Kruskal uniqueness prove that the solutions of (55) return exactly \(\{[c_j]:j\in A_\nu\}\). The Veronese map has no extra real projective collision: if \(z^{\otimes s}\) and \(w^{\otimes s}\) are proportional for \(s\ge1\) while \(z,w\) are not, choose a linear functional vanishing on \(w\) but not on \(z\) and contract all \(s\) modes, a contradiction.

By the established Marcinkiewicz cumulant-tail lemma, each nondegenerate non-Gaussian source with all moments has infinitely many nonzero cumulants of order at least three. Consequently, for every source \(j\), there is some \(\nu_j\ge2n-1\) with \(\lambda_{j,\nu_j}\ne0\). Since there are finitely many sources,
\[
 K_{C,\varepsilon}=\max_{1\le j\le n}\nu_j<\infty. \tag{56}
\]
Starting at \(2n-1\), solve (55) and take the set-theoretic union of its recovered projective directions. Pairwise nonproportionality prevents merging, and the rule stops when the union reaches its known cardinality \(n\), no later than (56). Labels need not be transported across orders. Ratios and deletion ranks are unchanged by nonzero column scaling and permutation, so the vector frontier and every scalar projection are then computed from the recovered projective columns. The sharp-effect-frontier theorem proves that these observable maps equal the full observational-law-fiber causal targets.

We now prove that neither the stopping order nor inference can be uniform under only the qualitative source assumptions. Fix any finite \(K\ge3\) and \(0<\gamma<\pi/4\). The established finite-moment near-Gaussian witness supplies a centered variance-one, non-Gaussian, moment-determinate law \(F\) whose first \(K\) moments equal the Gaussian moments. Give two independent sources law \(F\), and use
\[
 C_0=I_2,
 \qquad
 C_1=\begin{pmatrix}\cos\gamma&-\sin\gamma\\
 \sin\gamma&\cos\gamma\end{pmatrix}. \tag{57}
\]
The two matrices have normalized, nonzero, pairwise nonproportional columns. Independence makes every joint source moment through total degree \(K\) equal its Gaussian counterpart, and orthogonal invariance of the standard bivariate Gaussian then makes the observed moments, hence cumulants, through order \(K\) identical under \(C_0\) and \(C_1\). Their projective direction sets differ by the fixed angle \(\gamma\). Therefore no inverse with a uniformly fixed finite maximum order can recover every member of the qualitative class.

For the sampling obstruction, for each \(N\) apply the near-Gaussian witness with any order at least three and with \(d_{\mathrm{TV}}(F_N,\Phi)\le1/(4N^2)\), where \(\Phi=\mathcal N(0,1)\). Let \(P_{0,N}\) and \(P_{1,N}\) be the observed laws under (57) with independent \(F_N\) sources. Total variation contracts under a measurable map and satisfies the two-factor telescoping bound. Both orthogonal images of \(\Phi^{\otimes2}\) equal \(\mathcal N(0,I_2)\), so
\[
 d_{\mathrm{TV}}(P_{0,N},P_{1,N})
 \le4d_{\mathrm{TV}}(F_N,\Phi)\le N^{-2}. \tag{58}
\]
The analogous \(N\)-factor telescoping bound yields
\[
 d_{\mathrm{TV}}(P_{0,N}^{\otimes N},P_{1,N}^{\otimes N})
 \le N^{-1}\longrightarrow0. \tag{59}
\]
If an estimator of the direction set were uniformly consistent, assigning its output to the nearer of the two fixed projective sets in (57) would produce a test with both errors tending to zero. But for every test region \(E\), the definition of total variation gives
\[
 P_{0,N}^{\otimes N}(E^c)+P_{1,N}^{\otimes N}(E)
 \ge1-d_{\mathrm{TV}}(P_{0,N}^{\otimes N},P_{1,N}^{\otimes N})\longrightarrow1, \tag{60}
\]
a contradiction.

The causal response frontiers are separated too. For treatment \(x=1\), direct evaluation of the one-by-one deletion minors gives
\[
 \mathcal R_1^{\mathrm{vec}}(C_0)=\{0\},
 \qquad
 \mathcal R_1^{\mathrm{vec}}(C_1)=\{\tan\gamma,-\cot\gamma\}. \tag{61}
\]
Thus the same nearest-target test rules out uniform Hausdorff consistency for the frontier. More precisely, let \(\mathscr C_N\) be a random collection of compact frontier sets, let coverage mean that the true frontier belongs to \(\mathscr C_N\), and measure its uncertainty by \(\operatorname{diam}_H(\mathscr C_N)=\sup_{R,R'\in\mathscr C_N}d_H(R,R')\). If uniform coverage is \(1-\alpha\) with \(\alpha<1/2\), then under the first experiment \(\mathscr C_N\) contains the first frontier with probability at least \(1-\alpha\). By (59), it contains the second frontier under that same experiment with probability at least \(1-\alpha-o(1)\). Both events therefore occur with probability at least \(1-2\alpha-o(1)>0\). On their intersection, \(\operatorname{diam}_H(\mathscr C_N)\) is at least the fixed positive distance between the two sets in (61), contradicting uniformly vanishing Hausdorff uncertainty.

Finally, the causal map itself is discontinuous without denominator and deletion-rank margins. After nonzero column normalization, for
\[
 D_t=\begin{pmatrix}t&1\\1&1\end{pmatrix},\qquad t\to0,
\]
one has \(\mathcal R_{12}(D_t)=\{1/t,1\}\) for \(t\ne0,1\), whereas \(\mathcal R_{12}(D_0)=\{1\}\). For
\[
 E_0=\begin{pmatrix}1&1&0&1\\1&0&0&0\\0&0&1&1\end{pmatrix}
\]
and \(E_t\) obtained by replacing entry \((2,2)\) by \(t\ne0\), direct deletion-rank calculation gives \(\mathcal R_{12}(E_0)=\{0\}\) but \(\mathcal R_{12}(E_t)=\{0,t,1\}\). These matrices retain full row rank and nonzero pairwise nonproportional columns near zero. Since scalar frontiers are coordinate projections of vector frontiers, both examples also obstruct a uniform vector modulus. Quantitative lower cumulants, moments, Veronese conditioning, angular and pencil separation, deletion ranks, and denominators are therefore genuinely required for robust finite-order inference. The proved common-order theorem imposes precisely such margins and hence retains its simultaneous scalar/vector coverage and root-\(N\) Hausdorff rates.